\SourceLeaf{326}

\begin{center}
[III.]

\textbf{DAS VIERECK.}

[1.]

[\emph{Handbuch, den astronomischen Wissenschaften gewidmet}, S.\,150.]
\end{center}

$p, p', p'', p'''$ seien die komplexen Orte der Punkte $(0), (1), (2), (3)$. Man betrachte das Dreieck, dessen Ecken $(4), (5), (6)$ durch
\[
pp'+p''p''', \qquad pp''+p'''p', \qquad pp'''+p'p'',
\]
und dessen Seiten [65, 46, 54] nach Größe und Richtung durch die Differenzen
\[
(p-p')(p''-p'''), \qquad (p-p'')(p'''-p'), \qquad (p-p''')(p'-p'')
\]
bestimmt sind, also wie folgt

\begin{center}
\begin{tabular}{rllcl}
65: & Größe & $(01).(23)$ & Richtung & $01+23$\\
46: & ,, & $(02).(31)$ & ,, & $02+31$\\
54: & ,, & $(03).(12)$ & ,, & $03+12$,
\end{tabular}
\end{center}

Daraus ergeben sich die Gleichungen (siehe\,S.\,201 und S.\,269 [des Handbuchs])
\[
\begin{aligned}
\frac{\sin(02+31-03-12)}{(01).(23)} &= \frac{\sin(03+12-01-23)}{(02).(31)}\\
&= \frac{\sin(01+23-02-31)}{(03).(12)} = \Omega.
\end{aligned}
\]

Ferner ergibt die Entwicklung eines dieser Ausdrücke leicht (siehe\,S.\,284 [des Handbuchs])
\[
\Omega = \frac{(01)\sin(03-02)+(02)\sin(01-03)+(03)\sin(02-01)}{(12).(13).(23)}.
\]

Bezeichnet man $\Omega$, insofern es von den oben festgelegten Elementen abhängt, mit $\Omega^{0123}$, so lautet die Bedingungsgleichung zwischen den Differentialen, wenn $(1), (2), (3), (4)$ feste Punkte und $(0)$ einen beweglichen Punkt bezeichnen und auch die Orientierung als veränderlich betrachtet wird
\[
\left.
\begin{aligned}
&(01).(23).(24).(34).\Omega^{0234}\,d(01)\\
-\;&(02).(13).(14).(34).\Omega^{0134}\,d(02)\\
[+\;&(03).(12).(14).(24).\Omega^{0124}\,d(03)\\
-\;&(04).(12).(13).(23).\Omega^{0123}\,d(04)
\end{aligned}
\right\} = 0.
\]
